% !TeX root = ../thuthesis-example.tex
%
% 第二章：量子网络物理建模
%

\chapter{量子网络物理建模}
\label{chap:phymodel}

本章给出本文采用的多跳量子中继网络物理模型。该模型以路径为评价对象：一条候选路径由若干量子链路段组成，每一段链路均携带初始保真度、建链时间、噪声类型以及端点相干时间等参数。路由层在扩展候选路径时，根据完整路径段序列重新计算路径时间瓶颈 $T_{\max}$、各段等待时间、存储退相干后的 Bell 对角态以及递归纠缠交换后的端到端保真度。后续章节对路由度量非保序性、\texttt{DijkstraRoutingProtocol} 基线以及 \texttt{SlicedDijkstraRoutingProtocol} 的分析，均建立在本章定义的路径级物理语义之上。

\section{多跳量子中继网络的物理模型}
\label{sec:network-model}

\subsection{节点与链路参数}

考虑由若干量子中继节点构成的一般网络图。对一条候选路径

\[
p = (v_0, v_1, \ldots, v_h),
\]

其第 $j$ 段链路

\[
e_j = (v_{j-1}, v_j)
\]

记为一条路径段。本文在路由标签中保留各路径段的物理参数，用于后续的路径级保真度计算。每一段链路包含以下信息：

\begin{itemize}
  \item 左右端节点名；
  \item 初始保真度 \texttt{initialFidelity}；
  \item 噪声族 \texttt{noiseFamily}；
  \item 量子链路建立时间 \texttt{setupTimeMs}；
  \item 左右端节点的相干时间 $\tau$。
\end{itemize}

上述参数由链路层和物理层共同给出。\texttt{initialFidelity} 描述纠缠生成后单跳 Bell 对的初始质量，其数值受信道噪声和纠缠制备过程影响；\texttt{setupTimeMs} 表示该链路建立纠缠所需的时间，优先取 \texttt{quantumSetupTimeMs}，未显式给出时再依次由 \texttt{UseCurrentLinkDelay}、\texttt{latency} 与 \texttt{classicalControlDelayMs} 推定；$\tau$ 表示端点量子存储器的相干时间。链路可用性由 \texttt{isAvailable} 与 \texttt{successRate} 共同约束：当 \texttt{isAvailable} 为假或 \texttt{successRate} $\le 0$ 时，该链路不参与路径扩展；对于已经进入候选集合的链路，本文的 \texttt{predicted\_fidelity} 度量不再额外引入 \texttt{successRate} 因子。

对应地，单跳链路的初始 Bell 态仍按 Werner 态参数化：
\begin{equation}
  \rho_i^{(0)}
  = F_i^{(0)} \vert\Phi^+\rangle\langle\Phi^+\vert
  + \frac{1 - F_i^{(0)}}{3}
    \left(
      \vert\Phi^-\rangle\langle\Phi^-\vert +
      \vert\Psi^+\rangle\langle\Psi^+\vert +
      \vert\Psi^-\rangle\langle\Psi^-\vert
    \right)
  \label{eq:werner-init}
\end{equation}
其中$\vert\Phi^\pm\rangle,\vert\Psi^\pm\rangle$为四个Bell基态。

\subsection{纠缠交换协议的时序模型}

本文采用并行建链、统一对齐的时序假设。对路径 $p$ 上的各段链路，纠缠生成在同一轮路径评估中并行开始；随后以最晚完成的链路作为交换前的对齐时刻。具体过程如下：

\begin{enumerate}
  \item 路径上的各段链路在同一轮评估中同时开始建立纠缠；
  \item 每段链路的就绪时间由其 \texttt{setupTimeMs} 给出；
  \item 整条路径的统一对齐时刻定义为
  \begin{equation}
    T_{\max}(p) = \max_{1 \le j \le h} T_j,
    \qquad T_j \equiv \texttt{setupTimeMs}(e_j);
    \label{eq:path-tmax}
  \end{equation}
  \item 更快完成的链路必须等待较慢链路，于是第 $j$ 段的等待时间为
  \begin{equation}
    W_j(p) = T_{\max}(p) - T_j;
    \label{eq:path-wait}
  \end{equation}
  \item 交换操作在路由度量中视为瞬时完成，不额外计入建链时间。
\end{enumerate}

在该时序模型下，存储退相干由路径内部不同链路的建链时间差引入。较早完成的链路在等待对齐时发生退相干，等待后的单跳 Bell 对再通过纠缠交换递归合成为端到端 Bell 对。因此，$T_{\max}$ 同时决定路径的时间瓶颈和各段链路的等待长度，也是后续切片路由算法选择 $T_{\max}$ 作为分组变量的物理依据。

\section{端到端纠缠保真度的解析公式}
\label{sec:fidelity-formula}

\subsection{退相干模型}

本文采用的存储退相干模型作用于单个存储量子比特。设量子比特在存储器中等待时间为 $\Delta t$，节点相干时间为 $\tau$，则等待过程由相位翻转信道表示为
\begin{equation}
  \mathcal{E}_{\lambda}(\rho)
  = \frac{1 + \lambda}{2}\rho
  + \frac{1 - \lambda}{2}(Z\rho Z),
  \qquad
  \lambda(\Delta t, \tau) = e^{-\Delta t / \tau},
  \label{eq:phase-flip-memory}
\end{equation}
对一对以式~\eqref{eq:werner-init}初始化的 Bell 对，若左右两端量子比特分别等待 $\Delta t_L,\Delta t_R$，则等待后的二体量子态仍保持 Bell 对角形式：
\begin{equation}
  \rho_i(\Delta t_L, \Delta t_R)
  = a_i \vert\Phi^+\rangle\langle\Phi^+\vert
  + b_i \vert\Phi^-\rangle\langle\Phi^-\vert
  + c_i \vert\Psi^+\rangle\langle\Psi^+\vert
  + d_i \vert\Psi^-\rangle\langle\Psi^-\vert,
  \label{eq:single-link-bell-diagonal}
\end{equation}
其中
\begin{align}
  q_i &\triangleq \frac{1 - F_i^{(0)}}{3}, \notag\\
  \Lambda_i &\triangleq
    \exp\!\left(-\frac{\Delta t_L}{\tau_i} - \frac{\Delta t_R}{\tau_{i+1}}\right), \notag\\
  a_i &= \frac{2F_i^{(0)} + 1 + \Lambda_i(4F_i^{(0)} - 1)}{6}, \notag\\
  b_i &= \frac{2F_i^{(0)} + 1 - \Lambda_i(4F_i^{(0)} - 1)}{6}, \notag\\
  c_i &= d_i = q_i. \label{eq:single-link-bell-weights}
\end{align}
因此，该链路在等待后的 Bell 保真度为
\begin{equation}
  F_i(\Delta t_L, \Delta t_R)
  = \langle\Phi^+\vert \rho_i(\Delta t_L, \Delta t_R) \vert\Phi^+\rangle
  = a_i,
  \label{eq:single-link-fidelity}
\end{equation}
式~\eqref{eq:single-link-fidelity} 表明，双端存储退相干对 Bell 保真度的影响由两端等待时间和相干时间共同决定，并通过 Bell 对角态权重重新分配体现。

\subsection{多跳端到端保真度}

设 Bell 对角态的 Bell 权重向量记为
\begin{equation}
  w = (a, b, c, d),
  \qquad a + b + c + d = 1,
\end{equation}
其中 $a$ 即相对于 $\vert\Phi^+\rangle$ 的 Bell 保真度。对两端等待导致的相位翻转退相干，可定义等待算子
\begin{equation}
  \mathcal{M}_{\Lambda}(a, b, c, d)
  =
  \left(
    \frac{1 + \Lambda}{2}a + \frac{1 - \Lambda}{2}b,\;
    \frac{1 - \Lambda}{2}a + \frac{1 + \Lambda}{2}b,\;
    \frac{1 + \Lambda}{2}c + \frac{1 - \Lambda}{2}d,\;
    \frac{1 - \Lambda}{2}c + \frac{1 + \Lambda}{2}d
  \right),
  \label{eq:bell-memory-operator}
\end{equation}
其中 $\Lambda$ 由式~\eqref{eq:single-link-bell-weights} 给出。

若两段相邻路径的端到端态分别为 $w=(a,b,c,d)$ 与 $w'=(a',b',c',d')$，则一次理想纠缠交换后的输出仍为 Bell 对角态，其 Bell 权重满足
\begin{equation}
  \mathcal{S}(w, w')
  =
  \bigl(
    aa' + bb' + cc' + dd',\;
    ab' + ba' + cd' + dc',\;
    ac' + ca' + bd' + db',\;
    ad' + da' + bc' + cb'
  \bigr).
  \label{eq:bell-swap-operator}
\end{equation}
由此，多跳端到端纠缠保真度可通过 Bell 对角态递推得到：首先根据式~\eqref{eq:path-wait} 对各跳 Bell 对施加等待算子 $\mathcal{M}_{\Lambda}$，随后沿路径顺序递归施加交换算子 $\mathcal{S}$。递推结束后，输出 Bell 对角态的第一个分量即为端到端 $\vert\Phi^+\rangle$ 保真度。

在 qns-3 的路由标签中，上述路径级计算结果分别记录为 \texttt{predicted\_fidelity} 与 \texttt{t\_max\_ms}。特别地，在\emph{无额外等待}且输入均为 Werner 态的情形，若两条相邻链路初始保真度分别为 $F_1,F_2$，则一次纠缠交换后的 Bell 保真度为
\begin{equation}
  F_{\mathrm{swap}}
  = F_1 F_2 + \frac{(1 - F_1)(1 - F_2)}{3}
  = \frac{4F_1F_2 - F_1 - F_2 + 1}{3},
  \label{eq:werner-swap-fidelity}
\end{equation}
式~\eqref{eq:single-link-fidelity} 和式~\eqref{eq:werner-swap-fidelity} 分别给出了存储等待与理想纠缠交换在本文模型下的闭式表达。二者共同构成后续路径级路由度量的物理基础。
